\section{Further questions and boundaries of the results}
The finite-order method leaves several natural questions. Each requires
additional arguments beyond the theorems proved here.

\paragraph{Effective inverse constants.}
The local analytic bounds and Chernoff inequalities are explicit in form,
but this article has not optimized their constants or implemented a certified
finite onset. A useful next step is an interval-arithmetic inverse with a
provable error radius, followed by exact count comparisons for any
ambiguous adjacent thresholds.

\paragraph{A better moderate size approximation.}
The $p=3$ diagnostic suggests retaining more of the defect in exponential
form rather than immediately collecting fractional powers. The full
$H$ truncation is already much better at one displayed input, but this
alone proves no generally superior scheme. Uniform relative error bounds
for an appropriately resummed approximation would be more informative
than choosing a truncation by appearance.

\paragraph{Critical marked laws.}
At $p=1$ the leading quadratic tilt is order one in the partition function.
The count proof resums it, but the sharp marked-law argument in
Section~\ref{spb:ep:sec:marked} relies on $\lambda^2/n\to0$.
A separate centered analysis is needed to determine the best conditional
reference and sharp distances at criticality. Nothing in the inverse
transfer proves those probability results.
\begin{writenote}[Answered in Part~II, 7 October 2026]
Part~II (Report~235) carries out this centered analysis for the inactive
count: the sharp distances to both conditioned Poisson references
(Corollary~\ref{spb:ml:cor:Poisson}, with the constants of~\eqref{spb:ep:eq:TVlambda}
and~\eqref{spb:ep:eq:TVmu}), all fixed-order marked cumulants
(Theorem~\ref{spb:ml:thm:cumulants}), and the globally optimal
parity-conditioned binomial with its sharp constant
(Theorem~\ref{spb:ml:thm:global}). It does not determine a best reference
outside the conditioned binomial family.
\end{writenote}

\paragraph{Residue effects beyond algebraic order.}
The polynomial root filter displays exponentially small terms at nonzero
roots of unity. A remainder that is merely algebraically small can swamp
those terms. Establishing a secondary sector of the full deformed sum
would require a direct complex or oscillatory remainder analysis with
relative accuracy on that smaller scale. No such claim follows from
\eqref{spb:ep:eq:filter} alone.

\paragraph{Large order and growing parameters.}
All estimates hold at every chosen fixed order, with constants depending on
that order and on $p$. They give neither convergence of the infinite
asymptotic series nor optimal truncation. The scale $c_p$ grows rapidly
with $p$, so estimates uniform for a growing $p$ would need a fresh analysis.
There is also no continuous transition theorem joining $p=1$ to $p\ge2$;
the family considered here has integer $p$.

\paragraph{Sharper marked approximations.}
The tilted reference removes the leading location perturbation and leaves
a second Charlier correction. Higher signed Charlier corrections are
available algebraically. Turning them into positive probability
approximants with explicit distance bounds, or proving finer moderate
and large deviation estimates, would extend the present central laws.

\paragraph{Broader graph statistics.}
The inactive count determines the numbers of active and total vertices
through affine relations, but does not describe components, cycles,
indegrees, or repeated targets. Those observables require additional marks
in the directed-graph model. The exact size statistic and label convention
should be retained in any such extension.

The results established here are therefore concrete finite-order count,
probability, and inverse theorems. Their strengths are exact reduction,
global domination, explicit rational generators, and controlled error
transfer. Their scope does not depend on an unsupported claim about
exponentially small sectors or historical novelty.
\begin{writenote}[Added 7 October 2026, under Vladimir's standing rule of 4 October 2026]
Every open question of this Part is in this section, with what is missing;
the critical marked law is answered in Part~II within the conditioned binomial
family (note above). The write adds from the non-claims interval-certified values for the
tables and diagnostics of Section~\ref{spb:ep:sec:numerics}, which are
floating computations. No claim of the
source was found false; the OEIS refinements' first corrections are refuted
(Remark~\ref{spb:ep:rem:oeis}).
\end{writenote}
